\documentclass[11pt]{article}

\usepackage[a4paper,margin=0.88in]{geometry}
\usepackage{amsmath,amssymb,mathtools,bm}
\usepackage[dvipsnames]{xcolor}
\usepackage{booktabs}
\usepackage{enumitem}
\usepackage[hidelinks,hypertexnames=false]{hyperref}
\usepackage[capitalise]{cleveref}

\setlength{\parindent}{0pt}
\setlength{\parskip}{5pt}
\setlist[itemize]{leftmargin=*,nosep}
\setlist[enumerate]{leftmargin=*,nosep}

\newcommand{\wt}{\widetilde}
\newcommand{\Original}[1]{\textcolor{BrickRed}{#1}}
\newcommand{\Corrected}[1]{\textcolor{MidnightBlue}{#1}}
\newcommand{\Issue}[1]{\par\smallskip\noindent\textbf{\textcolor{BrickRed}{What is wrong.}} #1\par}
\newcommand{\Fix}[1]{\par\smallskip\noindent\textbf{\textcolor{MidnightBlue}{Replacement.}} #1\par}
\newcommand{\Added}[1]{\par\smallskip\noindent\textbf{\textcolor{ForestGreen}{Added here.}} #1\par}
\newcommand{\Why}[1]{\par\smallskip\noindent\textbf{Why.} #1\par}
\newcommand{\file}[1]{\texttt{#1}}

\title{Annotated Audit of \texttt{Calculations.tex}\\
Shortcut-to-absorbing-target correction, with original-step replacements}
\author{}
\date{31 May 2026}

\begin{document}
\maketitle
\tableofcontents

\section{Provenance and how to read this version}

This is the referee-facing annotated version of Luca's original
\file{Calculations.tex}.  It is intentionally organised like a markup pass over
the original manuscript: each critical original step is shown in red, followed
by the correction in blue and the reason the replacement is necessary.

The original and audit files are archived beside this document:
\begin{center}
\scriptsize
\begin{tabular}{@{}p{0.14\linewidth}p{0.52\linewidth}p{0.28\linewidth}@{}}
\toprule
Role & File & Notes \\
\midrule
Original source & \file{source/original\_from\_outlook\_Calculations.tex} &
Outlook attachment, Luca, 28 May 2026 \\
Original PDF & \file{source/original\_from\_outlook\_Calculations.pdf} &
Outlook attachment, Luca, 28 May 2026 \\
Audit trail & \file{Calculations\_marked\_changes\_v3.tex} &
Issue-by-issue marked changes \\
Full correction & \file{Calculations\_corrected\_complete\_v3.tex} &
Expanded corrected derivation \\
This file & \file{luca\_calculations\_annotated\_audit.tex} &
Original-step correction map \\
\bottomrule
\end{tabular}
\end{center}

Line anchors below refer to the archived original source.  The main mathematical
fault begins in the shortcut section around original lines 383--464.

\section{Executive correction}

\fbox{\begin{minipage}{0.96\linewidth}
For the stochastic model in which a fraction
\[
  \lambda=\beta(1-q)
\]
of the lazy self-loop at shortcut source $u$ is redirected to the absorbing
target $v$, the shortcut is an extra killing channel in the transient state
space.  Therefore the corrected resolvent is
\[
  \boxed{
  \wt S_{n_0}(n,z)=\wt W_{n_0}(n,z)
  -\frac{z\lambda\,\wt W_{n_0}(u,z)\wt W_u(n,z)}
  {1+z\lambda\,\wt W_u(u,z)}.}
\]
The original denominator $1-z\lambda\wt W_u(u,z)$ and the original positive
propagator correction have the wrong sign for a shortcut into an absorbing
target.
\end{minipage}}

The corrected first-passage generating function is
\begin{equation}
  \boxed{
  \wt F_{n_0}(v,z)=\wt F^{(0)}_{n_0}(v,z)+
  \frac{z\lambda\,\wt W_{n_0}(u,z)\bigl[1-\wt F^{(0)}_u(v,z)\bigr]}
  {1+z\lambda\,\wt W_u(u,z)}.}
  \label{eq:annotated-main-F}
\end{equation}
This is positive relative to the no-shortcut hitting generating function for
$0\leq z<1$, as a shortcut into the target must be.

\section{What can be retained from the original}

The absorbing-ring preliminaries are mostly usable after local corrections.  In
particular, the no-shortcut first-passage expansion and the killed-ring
propagator strategy are the right starting point:
\[
  \mathcal F_{n_0}^{(0)}(m,t)
  =\sum_{\ell=1}^{N/2}f_\ell(n_0,m)\alpha_\ell^{t-1},
  \qquad
  \alpha_\ell=1-q+q\cos\frac{(2\ell-1)\pi}{N},
\]
with $N$ even stated before using the $N/2$ odd-sector form.

\Issue{The original absorbing-ring reduction contains a sign-bookkeeping
problem in the double-sum intermediate line, and the subsequent
initial-condition check attributes zero to individual sums that only cancel
correctly in combination.}

\Fix{Keep the final killed-ring propagator after spelling out the sign
convention and the orthogonality sums.  This is a local repair, not the main
failure of the manuscript.}

\Added{The full v3 derivation records the corrected finite-sum identities and
the corrected orthogonality check; this annotated version uses those as the
baseline and focuses on the shortcut section, where the conclusion changes.}

\section{Critical step 1: the defect formula is not modular after absorption}

\textbf{Original location:} lines 383--399 of \file{Calculations.tex}.

The original manuscript first writes a probability-conserving one-defect
formula and then replaces the free-ring propagator $\wt P$ by the killed
propagator $\wt W$:
\begingroup\color{BrickRed}
\begin{align}
\wt S_{n_0}(n,z)
&=\wt W_{n_0}(n,z)
+\frac{z\beta(1-q)\wt W_{n_0}(u,z)
\left[\wt W_u(n,z)-\wt W_v(n,z)\right]}
{1-z\beta(1-q)\left[\wt W_u(u,z)-\wt W_v(u,z)\right]}.
\tag{O1}
\end{align}
\endgroup
When $v$ is the absorbing site, the original then simplifies this to
\begingroup\color{BrickRed}
\begin{align}
\wt S_{n_0}(n,z)
=\wt W_{n_0}(n,z)
+\frac{z\beta(1-q)\wt W_{n_0}(u,z)\wt W_u(n,z)}
{1-z\beta(1-q)\wt W_u(u,z)}.
\tag{O2}
\end{align}
\endgroup

\Issue{This replacement changes the model.  In the probability-conserving ring,
a shortcut transfers mass from $u$ to another transient site.  If the endpoint
$v$ is absorbing, the same transfer removes mass from the transient state
space.  The perturbation is therefore negative killing at $u$, not a positive
rank-one return term.}

\Fix{State the transient matrix first.  If $P_0$ is the killed-ring transient
matrix and $\lambda=\beta(1-q)$ is the shortcut-to-target probability, then}
\begin{equation}
  \Corrected{P=P_0-\lambda |u\rangle\langle u|.}
  \label{eq:transient-killing}
\end{equation}
Therefore
\begin{align}
  (I-zP)^{-1}
  &=\bigl(I-zP_0+z\lambda |u\rangle\langle u|\bigr)^{-1} \notag\\
  &= (I-zP_0)^{-1}
  -\frac{z\lambda (I-zP_0)^{-1}|u\rangle\langle u|(I-zP_0)^{-1}}
  {1+z\lambda\langle u|(I-zP_0)^{-1}|u\rangle}.
\end{align}
Taking the $(n,n_0)$ element gives
\begingroup\color{MidnightBlue}
\begin{align}
\wt S_{n_0}(n,z)=\wt W_{n_0}(n,z)
-\frac{z\lambda\wt W_{n_0}(u,z)\wt W_u(n,z)}
{1+z\lambda\wt W_u(u,z)}.
\tag{C1}
\end{align}
\endgroup

\Why{The sign is forced by probability monotonicity.  Adding a shortcut into an
absorbing target cannot increase transient occupation.  The original formula
does exactly that, so it cannot describe this stochastic model.}

\section{Critical step 2: the first-passage correction has the wrong sign}

\textbf{Original location:} lines 401--414.

The original first-passage equation is
\begingroup\color{BrickRed}
\begin{align}
\wt F_{n_0}(v,z)=\wt{\mathcal F}_{n_0}(v,z)
-\frac{z\beta(1-q)\wt W_{n_0}(u,z)
\left[1-\wt{\mathcal F}_u(v,z)\right]}
{1-z\beta(1-q)\wt W_u(u,z)}.
\tag{O3}
\end{align}
\endgroup

\Issue{This makes the shortcut contribution negative.  For $0\leq z<1$, the
quantities $\wt W_{n_0}(u,z)$ and $1-\wt F_u^{(0)}(v,z)$ have non-negative
coefficients.  A target shortcut should increase, not decrease, the
first-passage generating function.}

\Fix{Summing the corrected transient propagator over nonabsorbing $n$ and
using
\[
  \wt R_a^{(0)}(z)=\frac{1-\wt F_a^{(0)}(v,z)}{1-z}
\]
gives the corrected equation \eqref{eq:annotated-main-F}.}

\Added{This is also the simplest diagnostic to show a reader where the original
argument turns around: the sign of the shortcut contribution is visibly wrong
before any Chebyshev algebra is used.}

\section{Critical step 3: the auxiliary denominator and Fig. 1 geometry change}

\textbf{Original location:} lines 411--436.

The original defines
\begingroup\color{BrickRed}
\begin{equation}
\wt H_{N,|u-v|}(z)
=\frac{\beta(1-q)}{q}
\frac{1}{1-z\beta(1-q)\wt W_u(u,z)}
\tag{O4}
\end{equation}
\endgroup
and in the symmetric case $|u-v|=N/2$ obtains
\begingroup\color{BrickRed}
\begin{equation}
  \wt H_{N,N/2}(z)
  =\frac{1}{\dfrac{q}{\beta(1-q)}-\phi(z)}.
\tag{O5}
\end{equation}
\endgroup

\Issue{Because the denominator inherited the wrong sign, the pole picture based
on a positive horizontal line intersecting $\phi(1/s)$ is the wrong picture for
the stochastic shortcut-to-target model.}

\Fix{With the corrected killing denominator, the same symbolic object becomes}
\begingroup\color{MidnightBlue}
\begin{equation}
  \wt H^{\rm corr}_{N,N/2}(z)
  =\frac{1}{\dfrac{q}{\beta(1-q)}+\phi(z)}.
\tag{C2}
\end{equation}
\endgroup
Equivalently, if one keeps the plotted curve $\phi(1/s)$, the horizontal line
is $-q/[\beta(1-q)]$, not $+q/[\beta(1-q)]$.

\Added{The figure caption should also correct the typo ``given the requirement
$q>1$''.  The stochastic model requires $0<q<1$, and a nonzero shortcut
probability proportional to $1-q$ requires $q<1$.}

\section{Critical step 4: replace the symmetric-case answer}

For $N=2L$, absorbing target $v=0$, shortcut source $u=L$, and starting
distance $\rho$ from $v$, set
\[
  a=\frac{q}{\beta(1-q)},
  \qquad
  y=1+\frac{1}{q}\left(\frac1z-1\right).
\]
The corrected closed form is
\begin{equation}
  \Corrected{
  \wt F_\rho(z)=
  \frac{aT_{L-\rho}(y)+U_{\rho-1}(y)+U_{L-\rho-1}(y)}
  {aT_L(y)+U_{L-1}(y)},\qquad U_{-1}=0.}
  \label{eq:corrected-symmetric-F}
\end{equation}

\Issue{The original long convolution formula downstream of $H$ is attached to
the wrong denominator.  It also contains several local issues: the baseline
term uses $f_\ell(n_0,u)$ where the target term should be
$f_\ell(n_0,v)$, and odd-sector terms involving $g^{(2)}_r$ are paired with
$\gamma_r$ in places where the odd-sector eigenvalue is $\alpha_r$.}

\Fix{Use \eqref{eq:corrected-symmetric-F} as the primary corrected result.
Then obtain time-domain coefficients by residues at the roots of the corrected
denominator}
\begin{equation}
  \Corrected{D(y)=aT_L(y)+U_{L-1}(y).}
  \label{eq:corrected-D}
\end{equation}
If $y_j$ are its roots and $s_j=1-q+qy_j$, then
\begin{equation}
  \Corrected{
  F_\rho(t)=\sum_{j=1}^{L} B_{\rho j}s_j^{t-1},
  \qquad
  B_{\rho j}=\frac{q\,N_\rho(y_j)}{D'(y_j)},}
  \label{eq:residue-expansion}
\end{equation}
where
\[
  N_\rho(y)=aT_{L-\rho}(y)+U_{\rho-1}(y)+U_{L-\rho-1}(y).
\]

\section{Critical step 5: replace the beta-threshold conclusion}

\textbf{Original location:} line 464.

The original conclusion states that the second-peak condition is
\[
\Original{
  \beta\leq \frac{2q}{(1-q)N}}
\]
and that the right tail changes between $s_1^t$ and $t\alpha_1^t$ depending
on this threshold.

\Issue{That threshold comes from the wrong-sign pole geometry.  The repeated
pole that produces $t\alpha_1^t$ is an artefact of the uncorrected denominator,
not a property of the corrected stochastic shortcut-to-target model.}

\Fix{The corrected pole equation is}
\begin{equation}
  \Corrected{
  aT_L(y_s)+U_{L-1}(y_s)=0,\qquad
  y_s=1+\frac{s-1}{q}.}
  \label{eq:corrected-pole-y}
\end{equation}
Writing $y_s=\cos\theta$ gives
\begin{equation}
  \Corrected{\tan(L\theta)=-a\sin\theta.}
  \label{eq:corrected-pole-theta}
\end{equation}
The largest positive root satisfies
\begin{equation}
  \Corrected{
  \gamma_1<s_1<\alpha_1,\qquad
  \gamma_1=1-q+q\cos\frac{2\pi}{N},\quad
  \alpha_1=1-q+q\cos\frac{\pi}{N}.}
  \label{eq:root-ordering}
\end{equation}
Thus the corrected right tail is
\begin{equation}
  \Corrected{
  F_\rho(t)\sim B_{\rho1}s_1^{t-1},\qquad t\to\infty,}
  \label{eq:corrected-tail}
\end{equation}
not generically $t\alpha_1^t$.

\Added{The clean decay-dependence statement is the weak-shortcut expansion}
\begin{equation}
  \Corrected{
  s_1=\alpha_1-\frac{2\beta(1-q)}{N}
  +O\!\left([\beta(1-q)]^2\right),}
  \label{eq:s1-expansion}
\end{equation}
so
\begin{equation}
  \Corrected{
  1-s_1=1-\alpha_1+\frac{2\beta(1-q)}{N}
  +O\!\left([\beta(1-q)]^2\right).}
  \label{eq:gap-expansion}
\end{equation}
For large $N$,
\begin{equation}
  \Corrected{
  1-s_1\sim \frac{q\pi^2}{2N^2}
  +\frac{2\beta(1-q)}{N}.}
  \label{eq:large-N-gap}
\end{equation}

\section{What this annotated version adds}

\begin{enumerate}
\item It pins down the stochastic interpretation before using a defect formula:
the shortcut endpoint is absorbing, so the transient perturbation is killing.
\item It replaces the original sign pattern in the propagator and first-passage
generating function.
\item It replaces the positive horizontal-line pole picture by the corrected
$a+\phi=0$ picture.
\item It gives the corrected symmetric generating function, pole equation, and
tail expansion in one place.
\item It separates two questions that the original merges: the asymptotic
right-tail decay and the finite-time visual existence of a second peak.
\end{enumerate}

\section{Short referee-facing conclusion}

The absorbing-ring part of the original calculation can be repaired locally.
The shortcut-to-target part cannot be repaired by small notation changes,
because the original formula uses the probability-conserving shortcut defect
after the shortcut endpoint has become absorbing.  Under the stated stochastic
model, the shortcut is an extra killing channel at $u$.  This flips the sign of
the rank-one correction and changes the denominator from
\[
  1-z\beta(1-q)\wt W_u(u,z)
  \quad\hbox{to}\quad
  1+z\beta(1-q)\wt W_u(u,z).
\]
Consequently the claimed threshold
\[
  \beta\leq \frac{2q}{(1-q)N}
\]
is not a valid analytical condition for the corrected model, and the corrected
long-time tail is controlled by the simple pole $s_1$ solving
$aT_L(y_s)+U_{L-1}(y_s)=0$.

\end{document}
